\subsection{Delegasi Akses dan Atenuasi \textit{Offline}}
\label{subsec:implementasi-atenuasi-offline-delegasi}

Implementasi pemberian akses awal, delegasi akses, dan atenuasi \textit{offline} pada subbab ini merupakan realisasi dari \texttt{R-TA-04} dan \texttt{R-TA-08} melalui komponen \texttt{I-TA-02} dan \texttt{I-TA-06}.

Mekanisme ini merupakan realisasi dari rancangan pada Subbab \ref{subsec:rancangan-alur-akses-dan-delegasi-macaroon}. Implementasi dibagi menjadi dua bagian utama. Bagian pertama adalah proses pemberian akses awal dari pasien kepada petugas rekam medis atau petugas administratif sebagai penerima akses awal. Bagian kedua adalah proses delegasi akses antar tenaga kesehatan dalam satu fasilitas pelayanan kesehatan yang dilakukan melalui atenuasi token Macaroon secara lokal.

\subsubsection{Delegasi Akses dari Pasien kepada Petugas Rekam Medis}
\label{subsubsec:delegasi-akses-dari-pasien-kepada-petugas-rekam-medis}

Implementasi delegasi akses dari pasien kepada petugas rekam medis merupakan realisasi dari rancangan pada Subbab \ref{subsubsec:penerbitan-token-awal-oleh-pasien-kepada-petugas-rekam-medis}.

Pada proses penerbitan token awal, pasien melakukan \textit{request} menggunakan \textit{Patient Client} untuk menerbitkan token awal kepada petugas rekam medis. Proses diawali dengan pasien melakukan \textit{scan} QR code dari \texttt{AdministrativePersonnel}. QR code tersebut memuat sebuah \textit{string} yang dipisahkan dengan tanda \texttt{@} dengan format \texttt{\{hospitalCid\}@\{personnelIota\}@\{personnelPrePK\}}. Data tersebut terdiri dari \texttt{hospital\_cid} sebagai ID identitas rumah sakit, \texttt{personnel\_iota\_address} sebagai alamat IOTA milik personal tenaga medis, dan \texttt{personnel\_pre\_public\_key} sebagai PRE public key personal tenaga medis dalam format base64 yang digunakan untuk enkripsi \texttt{data\_pre\_secret\_key\_seed}.

Setelah QR code berhasil di-\textit{scan}, \textit{Patient Client} mengambil informasi publik petugas dari \textit{smart contract} untuk memastikan identitas petugas yang akan menerima akses. Pasien kemudian memilih jenis layanan yang sedang berlangsung, yaitu \texttt{RAWAT\_JALAN} atau \texttt{RAWAT\_INAP}. Pilihan ini diperlukan karena menentukan cakupan dan durasi token Macaroons. Pada implementasi, akses \texttt{RAWAT\_JALAN} berlaku selama 1 hari, sedangkan akses \texttt{RAWAT\_INAP} berlaku selama 3 hari.

Selanjutnya, \textit{Patient Client} meminta \textit{nonce} kepada \textit{Server} PRE melalui \textit{endpoint} \texttt{/nonce}. \textit{Nonce} tersebut ditandatangani menggunakan kunci IOTA pasien untuk membuktikan bahwa permintaan \textit{grant} benar-benar berasal dari pasien pemilik alamat IOTA terkait. Bersamaan dengan itu, \textit{Patient Client} membuat \texttt{data\_pre\_secret\_key\_seed} baru, menurunkan \texttt{data\_pre\_public\_key}, dan mengenkripsi seed tersebut menggunakan \texttt{personnel\_pre\_public\_key}. \textit{Patient Client} juga membuat \textit{key fragment} \texttt{k\_frag} dari PRE \textit{secret key} pasien menuju \texttt{data\_pre\_public\_key}. Material PRE ini nantinya disimpan pada Redis PRE agar server dapat melakukan \textit{re-encryption} ketika petugas mengakses data pasien. \texttt{k\_frag} memiliki \textit{time to live} (TTL) sesuai dengan konteks pelayanan.

Permintaan penerbitan akses kemudian dikirim ke \textit{endpoint} \texttt{/keys} dengan \textit{payload} yang mencakup alamat IOTA pasien, alamat IOTA petugas, \texttt{hospital\_cid}, \texttt{encounter\_dataset}, tanda tangan pasien, \texttt{k\_frag}, \texttt{data\_pre\_public\_key}, \texttt{enc\_data\_pre\_secret\_key\_seed}, \texttt{data\_pre\_secret\_key\_seed\_capsule}, \texttt{patient\_pre\_public\_key}, dan \texttt{signer\_pre\_public\_key}. \textit{Server} PRE kemudian memverifikasi tanda tangan pasien terhadap \textit{nonce}, memeriksa role petugas pada \textit{smart contract}, dan memastikan bahwa petugas yang di-\textit{scan} merupakan \texttt{AdministrativePersonnel}.

Jika validasi berhasil, \textit{Server} PRE membentuk \texttt{MacaroonKey} dari \texttt{RootKey} dan menerbitkan dua Macaroon awal untuk petugas administratif, yaitu token \texttt{READ} dan token \texttt{UPDATE}. Token \texttt{READ} memiliki \textit{caveat} utama berupa \texttt{patient\_address}, \texttt{root\_subject}, \texttt{read\_dataset\_in}, \texttt{read\_function\_in}, \texttt{expires\_before}, \texttt{max\_delegation\_depth}, \texttt{hospital\_cid}, dan \texttt{purpose = READ}. Token \texttt{UPDATE} memiliki \textit{caveat} berupa \texttt{patient\_address}, \texttt{related\_rme\_id}, \texttt{root\_subject}, \texttt{write\_dataset\_in}, \texttt{write\_function\_in}, \texttt{expires\_before}, \texttt{max\_delegation\_depth}, \texttt{hospital\_cid}, dan \texttt{purpose = UPDATE}.

Nilai \textit{caveat} \texttt{read\_dataset\_in} dan \texttt{read\_function\_in} untuk token \texttt{READ} adalah semua \textit{dataset category} dan \textit{function category}. Hal ini dikarenakan tenaga kesehatan membutuhkan akses keseluruhan data RME pasien untuk melakukan pelayanan klinis.

Untuk token \texttt{UPDATE}, \textit{Server} PRE juga membuat \texttt{related\_rme\_id} baru sebagai identitas episode rekam medis yang sedang dibuat atau diperbarui. ID tersebut menggunakan format \texttt{RME-\{YYYY\}-\{sequence\_6\_digit\}}. Sedangkan untuk token \texttt{READ} tidak membawa \textit{caveat} \texttt{related\_rme\_id}. Token \texttt{UPDATE} membawa \textit{caveat} \texttt{related\_rme\_id} untuk mengikat operasi \textit{update} pada episode RME saat episode pelayanan yang berlangsung saja. Nilai \textit{sequence} diperoleh dari \textit{counter} global per tahun, sehingga pasien yang berbeda tidak menghasilkan ID RME yang sama.

Setelah kedua token diterbitkan, \textit{Server} PRE menyimpan material akses PRE ke Redis dengan TTL sesuai jenis layanan. Entri di PRE memiliki \textit{key} \verb|{personnel_iota_address}@{patient_iota_address}|. \textit{Server} kemudian menghitung \textit{hash} dari token \texttt{READ} dan token \texttt{UPDATE}, lalu mengembalikan \texttt{access\_token\_read}, \texttt{access\_token\_update}, \textit{hash} masing-masing token, serta \texttt{related\_rme\_id} kepada \textit{Patient Client}.

Pada tahap akhir, \textit{Patient Client} menenkripsi token dan metadata akses menggunakan \texttt{personnel\_pre\_public\_key}, lalu mencatat \textit{grant} akses ke IOTA melalui \textit{smart contract} fungsi \texttt{create\_access}. Metadata yang dicatat berisi dua entri terenkripsi, yaitu metadata untuk token \texttt{READ} dan metadata untuk token \texttt{UPDATE}. Dengan demikian, token Macaroon awal telah dirilis kepada \texttt{AdministrativePersonnel} secara terenkripsi, hash token telah dicatat untuk keperluan audit dan revokasi, dan material PRE telah tersedia di Redis untuk proses akses data berikutnya.

\subsubsection{Delegasi Akses antar Tenaga Kesehatan}
\label{subsubsec:delegasi-akses-antar-tenaga-kesehatan}


Implementasi delegasi akses antar tenaga kesehatan merupakan realisasi dari rancangan pada Subbab \ref{subsec:delegasi-dari-petugas-rekam-medis-kepada-perawat-dan-dokter} dan \ref{subsec:delegasi-dari-dokter-kepada-laboratorium-dan-apoteker}. Pada implementasi ini, Petugas Rekam Medis memiliki \textit{role} \texttt{AdministrativePersonnel}. Tenaga medis lain seperti perawat dan dokter memiliki \textit{role} \texttt{MedicalPersonnel} dengan \textit{sub-role} masing-masing, sesuai dengan penjelasan pada Subbab \ref{subsubsec:verifikasi-role-tenaga-medis-untuk-akses-dataset}.

Token yang didelegasikan bukan merupakan token awal baru yang diterbitkan ulang oleh \textit{Server} PRE, melainkan token turunan dari token Macaroon yang sudah dimiliki oleh \textit{delegator}. Aktor yang menerima token turunan atau didelegasikan akses adalah \textit{delegatee}. Proses delegasi token Macaroon dilakukan oleh \textit{delegator} dengan menambahkan \textit{caveat} baru pada token Macaroon. Macaroon turunan yang dibentuk dibungkus dalam metadata delegasi terenkripsi sebelum dicatat pada penyimpanan \textit{on-chain} agar hanya dapat dibuka oleh \textit{delegatee}.

Proses delegasi dilakukan melalui \textit{Hospital Client}. Ketika \textit{delegator} ingin memberikan akses kepada tenaga kesehatan lain, \textit{client} terlebih dahulu memuat material kriptografis milik \textit{delegator}, yaitu PRE \textit{secret key}, alamat IOTA, dan pasangan kunci IOTA. Di sisi lain, data publik milik \textit{delegatee} diperoleh dari \textit{input} delegasi, meliputi alamat IOTA \textit{delegatee} dan PRE \textit{public key} \textit{delegatee}. Data \textit{delegatee} tersebut dimuat dari IOTA setelah \textit{delegator} memilih salah satu kandidat \textit{delegatee} dari \textit{smart contract} fungsi \texttt{get\_delegatee\_candidates}. PRE \textit{public key} tersebut digunakan agar metadata delegasi dan material akses yang diperlukan hanya dapat dibuka oleh \textit{delegatee}.

Selain membentuk token turunan, \textit{Hospital Client} juga melakukan proses \textit{re-wrap} terhadap \verb|data_pre_secret_key_seed|. \textit{Seed} tersebut sebelumnya terenkripsi agar dapat dibuka oleh \textit{delegator}. Pada saat delegasi, \textit{delegator} mendekripsi \textit{seed} tersebut menggunakan PRE \textit{secret key} miliknya, kemudian mengenkripsinya ulang menggunakan PRE \textit{public key} milik \textit{delegatee}. Dengan cara ini, \textit{delegatee} dapat memperoleh material yang diperlukan untuk membuka data yang didelegasikan, tanpa harus meminta pasien menerbitkan ulang akses dari awal.

Pembentukan token turunan dilakukan dengan memanggil fungsi \texttt{attenuate\_macaroon} dari \textit{crate} \texttt{decmed\-macaroon\-auth}. Fungsi tersebut menerima token induk dalam bentuk serialisasi Macaroon dan parameter atenuasi dalam bentuk \texttt{DelegationAttenuationParams}. Parameter tersebut mencakup alamat \textit{delegator}, alamat \textit{delegatee}, cakupan \textit{dataset} yang dapat dibaca atau ditulis, cakupan \textit{function category} yang dapat dibaca atau ditulis, batas waktu berlaku token, batas kedalaman delegasi, serta \texttt{related\_rme\_id} jika dan hanya jika belum ada. Token turunan kemudian dibentuk dengan menambahkan \textit{caveat} baru yang merepresentasikan batasan-batasan tersebut.

\textit{Caveat} yang ditambahkan pada token turunan mencakup \texttt{parent\_token\_hash}, \texttt{delegated\_by}, \texttt{delegated\_to}, cakupan akses, \texttt{expires\_before}, \texttt{max\_delegation\_depth}, dan \texttt{related\_rme\_id} apabila token berhubungan dengan episode RME tertentu. \texttt{parent\_token\_hash} berisi nilai hash dari token induk dan digunakan untuk mendukung mekanisme revokasi berantai. \textit{Caveat} \texttt{delegated\_by} dan \texttt{delegated\_to} mencatat relasi delegasi antar aktor. \textit{Caveat} cakupan akses, seperti \texttt{read\_dataset\_in}, \texttt{write\_dataset\_in}, \texttt{read\_function\_in}, dan \texttt{write\_function\_in}, digunakan untuk membatasi hak akses \textit{delegatee} agar tidak melebihi hak akses yang dimiliki \textit{delegator}.

Sebelum token turunan dianggap valid, parameter atenuasi divalidasi agar tidak memperluas hak akses token induk. Validasi tersebut memastikan bahwa waktu kedaluwarsa token turunan tidak lebih panjang daripada token induk, cakupan baca dan tulis token turunan merupakan subset dari cakupan token induk, dan nilai \texttt{max\_delegation\_depth} berkurang 1. Selain itu, \texttt{related\_rme\_id} hanya dapat ditambahkan apabila token induk belum memiliki \textit{caveat} tersebut. Dengan aturan ini, delegasi hanya dapat mempersempit hak akses, bukan memperluasnya.

Setelah token turunan dibentuk, \textit{Hospital Client} membuat bukti delegasi dalam bentuk \texttt{DelegationProofContext}. Bukti ini memuat \texttt{token\_hash}, alamat IOTA \textit{delegator}, alamat IOTA \textit{delegatee}, \texttt{related\_rme\_id}, waktu kedaluwarsa, dan \texttt{scope\_fingerprint}. Nilai \texttt{scope\_fingerprint} dihitung dari cakupan akses yang didelegasikan, sehingga perubahan terhadap cakupan akses setelah proses penandatanganan dapat terdeteksi. Konteks bukti delegasi kemudian diserialisasi dan ditandatangani menggunakan kunci IOTA \textit{delegator}. Tanda tangan ini berfungsi sebagai bukti bahwa delegasi benar-benar dibuat oleh aktor yang memiliki token induk.

Token turunan, \textit{hash} token, bukti delegasi, tanda tangan delegator, dan material PRE yang telah dienkripsi ulang kemudian dikemas sebagai metadata delegasi. Metadata tersebut dienkripsi menggunakan PRE \textit{public key} milik \textit{delegatee}. Dengan demikian, isi token dan material akses tidak disimpan dalam bentuk \textit{plaintext}, baik pada sisi \textit{client} maupun ketika metadata dicatat sebagai bagian dari transaksi delegasi.

Tahap berikutnya adalah pencatatan delegasi ke \textit{smart contract} melalui fungsi \texttt{create\_delegated\_access}. Pada tahap ini, \textit{smart contract} melakukan validasi bahwa delegator dan \textit{delegatee} merupakan tenaga kesehatan yang terdaftar pada rumah sakit yang sama, keduanya berada dalam status aktif, dan delegator memang memiliki akses yang masih berlaku terhadap pasien terkait. Jika validasi berhasil, \textit{smart contract} membuat data akses baru untuk \textit{delegatee} dengan mencatat alamat delegator pada atribut \texttt{delegated\_by} dan menaikkan nilai \texttt{delegation\_depth} berdasarkan kedalaman akses induk. Data akses tersebut kemudian dimasukkan ke tabel akses baca atau akses tulis milik \textit{delegatee} sesuai jenis akses yang didelegasikan.

Pencatatan \textit{on-chain} juga menghasilkan entri audit delegasi pada data pasien. Entri audit tersebut memuat informasi seperti jenis delegasi, waktu kejadian, alamat \textit{delegator}, subjek awal dari rantai delegasi, alamat \textit{delegatee}, jenis akses, \texttt{related\_rme\_id}, kedalaman delegasi, \textit{hash} token turunan, \textit{hash} token induk, dan waktu kedaluwarsa. Dengan adanya audit ini, riwayat perpindahan akses antar tenaga kesehatan dapat ditelusuri tanpa harus membuka isi token Macaroon yang sebenarnya.

Pada sisi Redis, proses delegasi tidak membuat penyimpanan material PRE baru yang berdiri sendiri untuk setiap relasi delegasi. Redis tetap digunakan untuk menyimpan material akses PRE dari pemberian akses awal. Ketika \textit{delegatee} mengakses data, \textit{Server} PRE membaca rantai delegasi dari token Macaroon untuk menentukan kandidat subjek yang dapat digunakan dalam pencarian material akses. Pencarian dilakukan berdasarkan \texttt{active\_subject}, \textit{delegator} dalam rantai delegasi, hingga \texttt{root\_subject}. Mekanisme ini diperlukan karena material PRE awal berasal dari pemberian akses pasien kepada penerima awal, sedangkan token turunan hanya membawa pembatasan akses dan metadata delegasi.

Setiap permintaan akses dari \textit{delegatee} tetap melewati \textit{middleware} otorisasi pada \textit{Server} PRE. Middleware melakukan deserialisasi token Macaroon dari \textit{Bearer token}, memverifikasi signature Macaroon menggunakan \texttt{RootKey}, menghitung \textit{effective capability} dari seluruh \textit{caveat}, membentuk rantai delegasi, serta memeriksa status revokasi pada Redis. Pemeriksaan revokasi dilakukan terhadap \textit{hash} token aktif, \textit{hash} token induk, revokasi akar, dan revokasi setiap relasi delegasi dalam rantai. Apabila salah satu entri revokasi ditemukan, maka token dianggap tidak valid meskipun signature Macaroon masih benar secara kriptografis.

Dengan mekanisme tersebut, delegasi akses antar tenaga kesehatan dapat dilakukan tanpa keterlibatan langsung pasien pada setiap perpindahan akses, tetapi tetap berada dalam batasan akses yang sebelumnya telah diberikan oleh pasien. Atenuasi lokal pada Macaroon memastikan bahwa delegatee tidak dapat memperoleh hak akses yang lebih luas daripada delegator. Sementara itu, pencatatan on-chain dan pemeriksaan revokasi di Redis memastikan bahwa proses delegasi tetap dapat diaudit dan dibatalkan apabila diperlukan.
